% Generated by `nl-assemble 1` from book/pipeline/work/ch001/scenes/*.final.nl.md.
% Do not edit by hand: fix the scene files and re-run. English source kept as `% EN:` lines.
\renewcommand{\nlchapterword}{EEN}

% EN: \chapter{A Day of Very Low Probability}
\chapter{Een dag van zeer lage waarschijnlijkheid}

\begin{chapterOpeningQuote}
\noindent
% EN: Beneath the moonlight glints a tiny fragment of silver, a fraction of a line…
Onder het maanlicht glimt een heel klein stukje zilver, een fractie van een lijn…

\vspace*{2ex}
% EN: (black robes, falling)
(zwarte \gl{gewaden}{robes}, vallend)

\vspace*{2ex}
% EN: …blood spills out in litres, and someone screams a word.
…liters bloed stromen naar buiten, en iemand schreeuwt een woord.
\end{chapterOpeningQuote}

% EN: \lettrine{E}{very} inch of wall space is covered by a bookcase. Each bookcase has six shelves, going almost to the ceiling. Some bookshelves are stacked to the brim with hardback books: science, maths, history, and everything else. Other shelves have two layers of paperback science fiction, with the back layer of books propped up on old tissue boxes\authorsnotefootnotemark{} or lengths of wood, so that you can see the back layer of books above the books in front. And it still isn’t enough. Books are overflowing onto the tables and the sofas and making little heaps under the windows.
\lettrine{T}{egen} elke centimeter muur staat een boekenkast. Elke boekenkast heeft zes planken, tot bijna aan het plafond. Sommige boekenplanken staan helemaal vol boeken met een harde kaft: wetenschap, wiskunde, geschiedenis en al het andere. Andere planken hebben twee rijen sciencefiction met een zachte kaft. De achterste rij staat wat hoger, op oude tissuedozen\authorsnotefootnotemark{} of op stukken hout. Zo kun je de achterste boeken boven de voorste boeken zien. En het is nog steeds niet genoeg. De boeken stromen over op de tafels en de banken en liggen in kleine stapels onder de ramen.
\authorsnotetext{I do this in my own home.}

% EN: This is the living-room of the house occupied by the eminent Professor Michael Verres-Evans, and his wife, Mrs~Petunia Evans-Verres, and their adopted son, Harry James Potter-Evans-Verres.
Dit is de woonkamer in het huis van de \gl{eminente}{eminent} professor Michael Verres-Evans, en zijn vrouw, mevrouw~Petunia Evans-Verres, en hun adoptiezoon, Harry James Potter-Evans-Verres.

% EN: There is a letter lying on the living-room table, and an unstamped envelope of yellowish parchment, addressed to \emph{Mr~H.~Potter} in emerald-green ink.
Er ligt een brief op de tafel in de woonkamer, en een envelop van geelachtig \gl{perkament}{parchment}, zonder postzegel. Op de envelop staat in \gl{smaragdgroene}{emerald-green} inkt: \emph{de heer~H.~Potter}.

% EN: The Professor and his wife are speaking sharply at each other, but they are not shouting. The Professor considers shouting to be uncivilised.
De Professor en zijn vrouw praten op scherpe toon tegen elkaar, maar ze schreeuwen niet. De Professor vindt schreeuwen onbeschaafd.

% EN: “You’re joking,” Michael said to Petunia. His tone indicated that he was very much afraid that she was serious.
“Je maakt een grapje,” zei Michael tegen Petunia. Zijn toon zei: hij was heel bang. Misschien was ze wel serieus.

% EN: “My sister was a witch,” Petunia repeated. She looked frightened, but stood her ground. “Her husband was a wizard.”
“Mijn zus was een heks,” herhaalde Petunia. Ze leek bang, maar ze gaf niet toe. “Haar man was een tovenaar.”

% EN: “This is absurd!” Michael said sharply. “They were at our wedding—they visited for Christmas—”
“Dit is absurd!” zei Michael scherp. “Ze waren op onze bruiloft—ze kwamen met Kerstmis op bezoek—”

% EN: “I told them you weren’t to know,” Petunia whispered. “But it’s true. I’ve seen things—”
“Ik heb tegen ze gezegd: jij mocht het niet weten,” fluisterde Petunia. “Maar het is waar. Ik heb dingen gezien—”

% EN: The Professor rolled his eyes. “Dear, I understand that you’re not familiar with the sceptical literature. You may not realise how easy it is for a trained magician to fake the seemingly impossible. Remember how I taught Harry to bend spoons? If it seemed like they could always guess what you were thinking, that’s called cold reading—”
De Professor rolde met zijn ogen. “Lieverd, ik begrijp het: jij kent de sceptische literatuur niet. Dit weet je misschien niet: een getrainde \gl{goochelaar}{stage magician} doet heel makkelijk iets onmogelijks na. Het lijkt alleen maar onmogelijk. Weet je nog? Ik leerde Harry lepels buigen. Misschien leek het zo: ze konden altijd je gedachten raden. Dat heet \gl{cold reading}{guessing from small cues}—”

% EN: “It wasn’t bending spoons—”
“Het was geen lepels buigen—”

% EN: “What was it, then?”
“Wat was het dan?”

% EN: Petunia bit her lip. “I can’t just tell you. You’ll think I’m—” She swallowed. “Listen. Michael. I wasn’t—always like this—” She gestured at herself, as though to indicate her lithe form. “Lily did this. Because I—because I \emph{begged} her. For years, I begged her. Lily had \emph{always} been prettier than me, and I’d…been mean to her, because of that, and then she got \emph{magic}, can you imagine how I felt? And I \emph{begged} her to use some of that magic on me so that I could be pretty too, even if I couldn’t have her magic, at least I could be pretty.”
Petunia beet op haar lip. “Ik kan het je niet zomaar vertellen. Dan vind je me—” Ze slikte. “Luister. Michael. Ik was niet—altijd zo—” Ze wees naar zichzelf, naar haar slanke lichaam, zo leek het. “Lily heeft dit gedaan. Omdat ik—omdat ik haar \emph{smeekte}. Jarenlang smeekte ik haar. Lily was \emph{altijd} al mooier dan ik, en daarom was ik…gemeen tegen haar. En toen kreeg zij \emph{magie}. Hoe voelde ik me toen, denk je? En ik \emph{smeekte} haar: gebruik een beetje van die magie op mij. Dan ben ik ook mooi. Haar magie kon ik niet krijgen, maar ik kon tenminste mooi zijn.”

% EN: Tears were gathering in Petunia’s eyes.
Er kwamen tranen in Petunia’s ogen.

% EN: “And Lily would tell me no, and make up the most ridiculous excuses, like the world would end if she were nice to her sister, or a centaur told her not to—the most ridiculous things, and I hated her for it. And when I had just graduated from university, I was going out with this boy, Vernon Dursley, he was fat and he was the only boy who would talk to me. And he said he wanted children, and that his first son would be named Dudley. And I thought to myself, \emph{what kind of parent names their child Dudley Dursley?} It was like I saw my whole future life stretching out in front of me, and I couldn’t stand it. And I wrote to my sister and told her that if she didn’t help me I’d rather just—”
“En Lily zei altijd nee. Ze bedacht de belachelijkste \gl{smoesjes}{excuses}: de wereld verging als ze aardig was voor haar zus. Of het mocht niet van een centaur—de belachelijkste dingen, en daarom haatte ik haar. En toen ik net klaar was met de universiteit, had ik een vriendje, Vernon Duffeling. Hij was dik, en verder wilde geen enkele jongen met me praten. En hij zei: hij wilde kinderen, en zijn eerste zoon moest Dudley heten. En ik dacht bij mezelf: \emph{wat voor ouder noemt zijn kind Dudley Duffeling?} Ik zag mijn hele verdere leven al voor me, en ik hield het niet uit. En ik schreef mijn zus een brief: als ze me niet hielp, dan wilde ik liever gewoon—”

% EN: Petunia stopped.
Petunia stopte.

% EN: “Anyway,” Petunia said, her voice small, “she gave in. She told me it was dangerous, and I said I didn’t care any more, and I drank this potion and I was sick for weeks, but when I got better my skin cleared up and I finally filled out and…I was beautiful, people were \emph{nice} to me,” her voice broke, “and after that I couldn’t hate my sister any more, especially when I learned what her magic brought her in the end—”
“Hoe dan ook,” zei Petunia met een klein stemmetje, “ze gaf toe. Ze zei: het is gevaarlijk. En ik zei: het maakt me niet meer uit. En ik dronk een toverdrank en ik was wekenlang ziek. Maar toen ik beter werd, werd mijn huid gaaf. En ik werd eindelijk wat voller en…ik was mooi. Mensen waren \emph{aardig} tegen me,” haar stem brak, “en daarna kon ik mijn zus niet meer haten. Vooral niet toen ik het hoorde: haar magie bracht haar uiteindelijk—”

% EN: “Darling,” Michael said gently, “you got sick, you gained some weight while resting in bed, and your skin cleared up on its own. Or being sick made you change your diet—”
“Schat,” zei Michael zacht, “je werd ziek. Je lag in bed en werd wat zwaarder. En je huid werd vanzelf gaaf. Of door je ziekte ging je anders eten—”

% EN: “She was a witch,” Petunia repeated. “I saw it.”
“Ze was een heks,” herhaalde Petunia. “Ik heb het gezien.”

% EN: “Petunia,” Michael said. The annoyance was creeping into his voice. “You \emph{know} that can’t be true. Do I really have to explain why?”
“Petunia,” zei Michael. Er kroop irritatie in zijn stem. “Dat kan niet waar zijn. Dat \emph{weet} je. Moet ik echt uitleggen waarom?”

% EN: Petunia wrung her hands. She seemed to be on the verge of tears. “My love, I know I can’t win arguments with you, but please, you have to trust me on this—”
Petunia \gl{wrong}{wrung} haar handen. Ze was bijna in tranen. “Mijn lief, ik weet het: van jou kan ik geen discussie winnen. Maar alsjeblieft, op dit punt moet je me vertrouwen—”

% EN: “\emph{Dad! Mum!}”
“\emph{Pap! Mam!}”

% EN: The two of them stopped and looked at Harry as though they’d forgotten there was a third person in the room.
Ze stopten allebei en keken naar Harry. Ze vergaten het even, zo leek het: er was nog een derde persoon in de kamer.

% EN: Harry took a deep breath. “Mum, \emph{your} parents didn’t have magic, did they?”
Harry haalde diep adem. “Mam, \emph{jouw} ouders hadden geen magie, toch?”

% EN: “No,” Petunia said, looking puzzled.
“Nee,” zei Petunia. Ze keek verbaasd.

% EN: “Then no-one in your family knew about magic when Lily got her letter. How did \emph{they} get convinced?”
“Dan wist niemand in je familie van magie toen Lily haar brief kreeg. Hoe raakten \emph{zij} overtuigd?”

% EN: “Ah…” Petunia said. “They didn’t just send a letter. They sent a professor from Hogwarts. He—” Petunia’s eyes flicked to Michael. “He showed us some magic.”
“Ah…” zei Petunia. “Ze stuurden niet alleen een brief. Ze stuurden een professor van Zweinstein. Hij—” Petunia’s ogen schoten even naar Michael. “Hij liet ons wat magie zien.”

% EN: “Then you don’t have to fight over this,” Harry said firmly. Hoping against hope that this time, just this once, they would listen to him. “If it’s true, we can just get a Hogwarts professor here and see the magic for ourselves, and Dad will admit that it’s true. And if not, then Mum will admit that it’s false. That’s what the experimental method is for, so that we don’t have to resolve things just by arguing.”
“Dan hoeven jullie geen ruzie te maken,” zei Harry met vaste stem. Hij hoopte, zonder veel hoop: misschien luisterden ze deze keer, deze ene keer, naar hem. “Als het waar is, halen we gewoon een professor van Zweinstein hierheen en zien we de magie met eigen ogen. En dan geeft pap toe: het is waar. En anders geeft mam toe: het is niet waar. Daar is de experimentele methode voor, zodat we dingen niet alleen met ruzie hoeven op te lossen.”

% EN: The Professor turned and looked down at him, dismissive as usual. “Oh, come now, Harry. Really, \emph{magic}? I thought \emph{you’d} know better than to take this seriously, son, even if you’re only ten. Magic is just about the most unscientific thing there is!”
De Professor draaide zich om en keek op hem neer, afwijzend als altijd. “Och, kom nou, Harry. Echt, \emph{magie}? Ik dacht: \emph{jij} neemt zoiets niet serieus, jongen, hoewel je pas tien bent. Je weet wel beter. Er is bijna niets zo onwetenschappelijk als magie!”

% EN: Harry’s mouth twisted bitterly. He was treated well, probably better than most genetic fathers treated their own children. Harry had been sent to the best primary schools—and when that didn’t work out, he was provided with tutors from the endless pool of starving students. Always Harry had been encouraged to study whatever caught his attention, bought all the books that caught his fancy, sponsored in whatever maths or science competitions he entered. He was given anything reasonable that he wanted, except, maybe, the slightest shred of respect. A Doctor teaching biochemistry at Oxford could hardly be expected to listen to the advice of a little boy. You would listen to Show Interest, of course; that’s what a Good Parent would do, and so, if you conceived of yourself as a Good Parent, you would do it. But take a ten-year-old \emph{seriously}? Hardly.
Harry trok een bitter gezicht. Zijn ouders behandelden hem goed, waarschijnlijk beter dan de meeste genetische vaders hun eigen kinderen. Harry ging eerder naar de beste basisscholen. Dat werkte niet. Toen kreeg hij eigen leraren uit de eindeloze voorraad arme studenten. Vond Harry iets interessant? Dan stimuleerden ze hem altijd: bestudeer het maar. Vond hij een boek leuk? Dan kochten ze het. Deed hij mee aan een wedstrijd in wiskunde of wetenschap? Dan betaalden ze alles. Hij kreeg alles, als het maar redelijk was. Behalve misschien een heel klein beetje respect. Een Doctor, docent biochemie in Oxford, luisterde niet naar het advies van een kleine jongen. Dat kon je ook moeilijk verwachten. Luisteren deed je natuurlijk wel: dat heette Belangstelling Tonen. Zo deed een Goede Ouder dat. En als je jezelf als een Goede Ouder zag, dan deed je dat dus ook. Maar een jongen van tien \emph{serieus} nemen? Nauwelijks.

% EN: Sometimes Harry wanted to scream at his father.
Soms wilde Harry tegen zijn vader schreeuwen.

% EN: “Mum,” Harry said. “If you want to win this argument with Dad, look in chapter two of the first book of the Feynman Lectures on Physics. There’s a quote there about how philosophers say a great deal about what science absolutely requires, and it is all wrong, because the only rule in science is that the final arbiter is observation—that you just have to look at the world and report what you see. Um…off the top of my head I can’t think of where to find something about how it’s an ideal of science to settle things by experiment instead of arguments—”
“Mam,” zei Harry. “Als je deze discussie met pap wilt winnen, kijk dan in hoofdstuk twee van het eerste deel van de Feynman Lectures on Physics. Daar staat een \gl{citaat}{quote}: filosofen zeggen heel veel over de absolute voorwaarden voor wetenschap, en het is allemaal fout. Want de wetenschap heeft maar één regel: de hoogste rechter is de \gl{waarneming}{observation}. Je moet gewoon naar de wereld kijken en dan zeggen: dit zie ik. Eh… er staat ergens iets over het ideaal van de wetenschap: je beslist dingen met een experiment, niet met een discussie. Maar waar staat dat? Dat weet ik zo snel niet—”

% EN: His mother looked down at him and smiled. “Thank you, Harry. But—” her head rose back up to stare at her husband. “I don’t want to win an argument with your father. I want my husband to, to listen to his wife who loves him, and trust her just this once—”
Zijn moeder keek omlaag naar hem en glimlachte. “Dank je, Harry. Maar—” haar hoofd kwam weer omhoog en ze staarde naar haar man. “Ik wil geen discussie met je vader winnen. Ik wil, ik wil dit van mijn man: hij moet naar zijn vrouw luisteren. Zij houdt van hem. En hij moet haar vertrouwen, deze ene keer—”

% EN: Harry closed his eyes briefly. \emph{Hopeless.} Both of his parents were just hopeless.
Harry deed even zijn ogen dicht. \emph{Hopeloos.} Zijn ouders waren allebei gewoon hopeloos.

% EN: Now his parents were getting into one of \emph{those} arguments again, one where his mother tried to make his father feel guilty, and his father tried to make his mother feel stupid.
Nu kregen zijn ouders weer \emph{zo'n} ruzie. Zijn moeder deed haar best: zijn vader moest zich schuldig voelen. Zijn vader deed zijn best: zijn moeder moest zich dom voelen.

% EN: “I’m going to go to my room,” Harry announced. His voice trembled a little. “Please try not to fight too much about this, Mum, Dad, we’ll know soon enough how it comes out, right?”
“Ik ga naar mijn kamer,” kondigde Harry aan. Zijn stem trilde een beetje. “Alsjeblieft, maak niet te veel ruzie, mam, pap. We weten het resultaat snel genoeg, toch?”

% EN: “Of course, Harry,” said his father, and his mother gave him a reassuring kiss, and then they went on fighting while Harry climbed the stairs to his bedroom.
“Natuurlijk, Harry,” zei zijn vader, en zijn moeder gaf hem een geruststellende kus. Daarna gingen ze door met ruziemaken, terwijl Harry de trap nam naar zijn slaapkamer.

% EN: He shut the door behind him and tried to think.
Hij deed de deur achter zich dicht en probeerde na te denken.

% EN: The funny thing was, he \emph{should} have agreed with Dad. No-one had ever seen any evidence of magic, and according to Mum, there was a whole magical world out there. How could anyone keep something like that a secret? More magic? That seemed like a rather suspicious sort of excuse.
Het gekke was: eigenlijk \emph{moest} hij pap gelijk geven. Niemand zag ooit bewijs van magie, en volgens mam was er ergens een hele wereld vol magie. Hoe kon iemand zoiets geheim houden? Nog meer magie? Dat leek een nogal verdacht soort smoes.

% EN: It should have been a clean case for Mum joking, lying or being insane, in ascending order of awfulness. If Mum had sent the letter herself, that would explain how it arrived at the letterbox without a stamp. A little insanity was far, far less improbable than the universe really working like that.
Het moest eigenlijk een duidelijk geval zijn: mam maakte een grapje, of ze loog, of ze was gek—op een schaal van erg naar erger. Stel: mam stuurde de brief zelf. Dan was er een verklaring: zo kwam de brief zonder postzegel in de brievenbus. Een beetje gek zijn was veel, veel minder \gl{onwaarschijnlijk}{improbable} dan dit: het universum werkt echt zo.

% EN: Except that some part of Harry was utterly convinced that magic was real, and had been since the instant he saw the putative letter from the Hogwarts School of Witchcraft and Wizardry.
Alleen wist een deel van Harry het helemaal zeker: magie was echt. En dat deel wist het al vanaf de eerste blik op de zogenaamde brief van Zweinsteins Hogeschool voor Hekserij en Hocus-Pocus.

% EN: Harry rubbed his forehead, grimacing. \emph{Don’t believe everything you think,} one of his books had said.
Harry wreef over zijn voorhoofd en trok een gezicht. \emph{Geloof niet al je gedachten,} zo stond het in een van zijn boeken.

% EN: But this bizarre certainty…Harry was finding himself just \emph{expecting} that, yes, a Hogwarts professor would show up and wave a wand and magic would come out. The strange certainty was making no effort to guard itself against falsification—wasn’t making excuses in advance for why there wouldn’t be a professor, or the professor would only be able to bend spoons.
Maar deze bizarre zekerheid…Harry merkte: hij \emph{verwachtte} het gewoon. Ja, er komt een professor van Zweinstein, hij zwaait met een \gl{toverstok}{wand}, en uit de toverstok komt magie. De vreemde zekerheid beschermde zichzelf niet tegen \gl{weerlegging}{falsification}. Die maakte vooraf geen smoesjes, zoals: ‘misschien komt er geen professor’, of ‘misschien kan de professor alleen maar lepels buigen’.

% EN: \emph{Where do you come from, strange little prediction?} Harry directed the thought at his brain. \emph{Why do I believe what I believe?}
\emph{Waar kom jij vandaan, vreemde kleine voorspelling?} Harry richtte de gedachte op zijn \gl{brein}{brain}. \emph{Waarom geloof ik wat ik geloof?}

% EN: Usually Harry was pretty good at answering that question, but in this particular case, he had no \emph{clue} what his brain was thinking.
Meestal kon Harry die vraag best goed beantwoorden, maar in dit specifieke geval had hij geen \emph{idee}: wat dacht zijn brein?

% EN: Harry mentally shrugged. A flat metal plate on a door affords pushing, and a handle on a door affords pulling, and the thing to do with a testable hypothesis is to go and test it.
Harry haalde in gedachten zijn schouders op. Een platte plaat van metaal op een deur nodigt uit tot duwen. Een handvat op een deur nodigt uit tot trekken. En met een testbare hypothese doe je één ding: die ga je testen.

% EN: He took a piece of lined paper from his desk, and started writing.
Hij pakte een vel papier met lijntjes van zijn bureau en begon met schrijven.

\begin{writtenNote}
% EN: \letterAddress{Dear Deputy Headmistress}
\letterAddress{Geachte onderdirectrice}
\end{writtenNote}

% EN: Harry paused, reflecting; then discarded the paper for another, tapping another millimetre of graphite from his mechanical pencil. This called for careful calligraphy.
Harry stopte even en dacht na. Toen gooide hij het vel weg en pakte een nieuw vel. Hij klikte nog een millimeter grafiet uit zijn \gl{vulpotlood}{mechanical pencil}. Dit vroeg om zorgvuldige \gl{kalligrafie}{calligraphy}.

\begin{writtenNote}
% EN: \letterAddress{Dear Deputy Headmistress Minerva McGonagall,}
\letterAddress{Geachte onderdirectrice Minerva Anderling,}

% EN: \letterAddress{Or Whomsoever It May Concern:}
\letterAddress{Of Wie Dit Verder Aangaat:}

% EN: I recently received your letter of acceptance to Hogwarts, addressed to Mr~H.~Potter. You may not be aware that my genetic parents, James Potter and Lily Potter (formerly Lily Evans) are dead. I was adopted by Lily’s sister, Petunia Evans-Verres, and her husband, Michael Verres-Evans.
Onlangs ontving ik uw brief aan de heer~H.~Potter over mijn \gl{toelating}{acceptance} tot Zweinstein. Misschien weet u dit niet: mijn genetische ouders, James Potter en Lily Potter (vroeger Lily Evans), zijn dood. Mijn adoptieouders zijn Lily’s zus, Petunia Evans-Verres, en haar man, Michael Verres-Evans.

% EN: I am extremely interested in attending Hogwarts, conditional on such a place actually existing. Only my mother Petunia says she knows about magic, and she can’t use it herself. My father is highly sceptical. I myself am uncertain. I also don’t know where to obtain any of the books or equipment listed in your acceptance letter.
Ik heb zeer veel belangstelling voor een plaats op Zweinstein, op voorwaarde dat die school echt bestaat. Alleen mijn moeder Petunia weet van magie, zegt ze. Zelf kan ze het niet gebruiken. Mijn vader is zeer sceptisch. Zelf twijfel ik. Ook over de lijst met boeken en materiaal in uw brief heb ik een vraag: waar kan ik die verkrijgen?

% EN: Mother mentioned that you sent a Hogwarts representative to Lily Potter (then Lily Evans) in order to demonstrate to her family that magic was real, and, I presume, help Lily obtain her school materials. If you could do this for my own family it would be extremely helpful.
Volgens mijn moeder stuurde u eerder een \gl{vertegenwoordiger}{representative} van Zweinstein naar Lily Potter (toen nog Lily Evans). Zo kon haar familie zien: magie is echt. En die persoon hielp Lily, neem ik aan, met haar spullen voor school. Zou u dit ook voor mijn eigen familie doen? Dat helpt ons zeer.

% EN: \letterClosing[Sincerely,]{Harry James Potter-Evans-Verres.}
\letterClosing[Hoogachtend,]{Harry James Potter-Evans-Verres.}
\end{writtenNote}

% EN: Harry added their current address, then folded up the letter and put it in an envelope, which he addressed to Hogwarts. Further consideration led him to obtain a candle and drip wax onto the flap of the envelope, into which, using a penknife’s tip, he impressed the initials H.J.P.E.V\@. If he was going to descend into this madness, he was going to do it with style.
Harry schreef hun huidige adres onder de brief. Hij vouwde de brief op en deed hem in een envelop. Op de envelop schreef hij het adres: Zweinstein. Na nog wat nadenken haalde hij een kaars en liet kaarsvet op de flap van de envelop vallen. Hij pakte een zakmes. Met de punt drukte hij zijn initialen in het kaarsvet: H.J.P.E.V\@. Als hij dan toch aan deze \gl{waanzin}{madness} begon, dan deed hij het met stijl.

% EN: Then he opened his door and went back downstairs. His father was sitting in the living-room and reading a book of higher maths to show how smart he was; and his mother was in the kitchen preparing one of his father’s favourite meals to show how loving she was. It didn’t look like they were talking to one another at all. As scary as arguments could be, \emph{not arguing} was somehow much worse.
Daarna deed hij zijn deur open en ging weer naar beneden. Zijn vader zat in de woonkamer en las een boek over hogere wiskunde. Zo liet hij zijn slimheid zien. Zijn moeder stond in de keuken en maakte een van zijn vaders favoriete maaltijden. Zo liet zij haar liefde zien. Blijkbaar praatten ze helemaal niet met elkaar. Ruzie kon eng zijn, maar \emph{geen ruzie maken} was op een of andere manier veel erger.

% EN: “Mum,” Harry said into the unnerving silence, “I’m going to test the hypothesis. According to your theory, how do I send an owl to Hogwarts?”
“Mam,” zei Harry in de akelige stilte, “ik ga de hypothese testen. Volgens jouw theorie: hoe stuur ik een uil naar Zweinstein?”

% EN: His mother turned from the kitchen sink to stare at him, looking shocked. “I—I don’t know, I think you just have to own a magic owl.”
Zijn moeder draaide zich om bij de \gl{gootsteen}{kitchen sink} en staarde hem geschokt aan. “Ik—ik weet het niet. Je moet gewoon een magische uil hebben, denk ik.”

% EN: That should’ve sounded highly suspicious, \emph{oh, so there’s no way to test your theory then}, but the peculiar certainty in Harry seemed willing to stick its neck out even further.
Dat moest eigenlijk heel verdacht klinken: \emph{o, dus je theorie kun je niet testen.} Maar de eigenaardige zekerheid in Harry wilde haar nek blijkbaar nog verder uitsteken.

% EN: “Well, the letter got here somehow,” Harry said, “so I’ll just wave it around outside and call ‘letter for Hogwarts!’ and see if an owl picks it up. Dad, do you want to come and watch?”
“Nou, de brief is hier op een of andere manier gekomen,” zei Harry, “dus ik ga gewoon buiten met de brief zwaaien en ‘brief voor Zweinstein!’ roepen. Dan zien we wel: pakt een uil hem? Pap, kom je kijken?”

% EN: His father shook his head minutely and kept on reading. \emph{Of course,} Harry thought to himself. Magic was a disgraceful thing that only stupid people believed in; if his father went so far as to \emph{test} the hypothesis, or even \emph{watch} it being tested, that would feel like \emph{associating} himself with that…
Zijn vader schudde heel licht zijn hoofd en las verder. \emph{Natuurlijk,} dacht Harry bij zichzelf. Magie was een schande. Alleen domme mensen geloofden in magie. Stel: zijn vader \emph{testte} zelfs de hypothese, of hij \emph{keek} alleen maar naar de test. Dan \emph{verbond} hij zichzelf met zoiets, zo voelde dat voor hem…

% EN: Only as Harry stumped out the back door, into the back garden, did it occur to him that if an owl \emph{did} come down and snatch the letter, he was going to have some trouble telling Dad about it.
Pas toen Harry door de achterdeur naar buiten stampte, de achtertuin in, bedacht hij iets. Als er \emph{wél} een uil naar beneden kwam en de brief greep, dan kon hij dat niet zo makkelijk aan pap vertellen.

% EN: \emph{But—well—that can’t \emph{really} happen, can it? No matter what my brain seems to believe. If an owl really comes down and grabs this envelope, I’m going to have worries a lot more important than what Dad thinks.}
\emph{Maar—nou ja—dat kan niet \emph{echt} gebeuren, toch? Mijn brein gelooft blijkbaar van wel—nou en. Als er echt een uil naar beneden komt en deze envelop pakt, dan heb ik veel belangrijkere zorgen dan de mening van pap.}

% EN: Harry took a deep breath, and raised the envelope into the air.
Harry haalde diep adem en stak de envelop in de lucht.

% EN: He swallowed.
Hij slikte.

% EN: Calling out \emph{Letter for Hogwarts!} while holding an envelope high in the air in the middle of your own back garden was…actually pretty embarrassing, now that he thought about it.
\emph{Brief voor Zweinstein!} roepen, met een envelop hoog in de lucht, midden in je eigen achtertuin. Dat was…eigenlijk best \gl{gênant}{embarrassing}, als hij het zo bekeek.

% EN: \emph{No. I’m better than Dad. I will use the scientific method even if it makes me feel stupid.}
\emph{Nee. Ik ben beter dan pap. Ik zal de wetenschappelijke methode gebruiken, ook als ik me dan dom voel.}

% EN: “Letter—” Harry said, but it actually came out as more of a whispered croak.
“Brief—” zei Harry, maar het klonk eigenlijk meer als een schor gefluister.

% EN: Harry steeled his will, and shouted into the empty sky, “\emph{Letter for Hogwarts! Can I get an owl?}”
Harry verzamelde al zijn wilskracht en schreeuwde de lege lucht in: “\emph{Brief voor Zweinstein! Mag ik een uil?}”

% EN: “Harry?” asked a bemused woman’s voice, one of the neighbours.
“Harry?” vroeg een verbaasde vrouwenstem, een van de buren.

% EN: Harry pulled down his hand like it was on fire and hid the envelope behind his back like it was drug money. His whole face was hot with shame.
Harry trok zijn hand snel omlaag, als van een hete kachel, en verstopte de envelop achter zijn rug als drugsgeld. Zijn hele gezicht was heet van schaamte.

% EN: An old woman’s face peered out from above the neighbouring fence, grizzled grey hair escaping from her hairnet. Mrs~Figg, the occasional babysitter. “What are you doing, Harry?”
Het gezicht van een oude vrouw verscheen boven het hek van de buren. Grijze plukken haar ontsnapten uit haar \gl{haarnetje}{hairnet}. Mevrouw~Figg, soms zijn oppas. “Wat ben je aan het doen, Harry?”

% EN: “Nothing,” Harry said in a strangled voice. “Just—testing a really silly theory—”
“Niets,” zei Harry met \gl{verstikte}{strangled} stem. “Gewoon—een heel domme theorie aan het testen—”

% EN: “Did you get your acceptance letter from Hogwarts?”
“Heb je je \gl{toelatingsbrief}{acceptance letter} van Zweinstein gekregen?”

% EN: Harry froze in place.
Harry bleef doodstil staan.

% EN: “Yes,” Harry’s lips said a little while later. “I got a letter from Hogwarts. They say they want my owl by the 31st of July, but—”
“Ja,” zeiden Harry’s lippen een poosje later. “Ik heb een brief van Zweinstein gekregen. Ze willen mijn uil vóór 31 juli, zeggen ze, maar—”

% EN: “But you don’t \emph{have} an owl. Poor dear! I can’t imagine \emph{what} someone must have been thinking, sending you just the standard letter.”
“Maar je \emph{hebt} geen uil. Och, arme jongen! Ik kan het me niet voorstellen. \emph{Wat} dachten ze toch? Ze stuurden jou alleen maar de gewone brief.”

% EN: A wrinkled arm stretched out over the fence, and opened an expectant hand. Hardly even thinking at this point, Harry gave over his envelope.
Een arm vol rimpels kwam over het hek, en een hand ging open en wachtte. Harry dacht nu nauwelijks nog na. Hij gaf zijn envelop af.

% EN: “Just leave it to me, dear,” said Mrs~Figg, “and in a jiffy or two I’ll have someone over.”
“Laat het maar aan mij over, lieverd,” zei mevrouw~Figg, “en \gl{in een wip}{in a jiffy} heb ik hier iemand.”

% EN: And her face disappeared from over the fence.
En haar gezicht verdween achter het hek.

% EN: There was a long silence in the garden.
Er was een lange stilte in de tuin.

% EN: Then a boy’s voice said, calmly and quietly, “What.”
Toen zei een jongensstem, kalm en zacht: “Wat.”
